\documentclass[10pt]{article}
\usepackage{amsmath,amsfonts,amssymb,mathtools,mathrsfs,bm,bbm,array,xcolor,geometry}
\allowdisplaybreaks
\setlength{\parindent}{0pt}
\setlength{\parskip}{4pt}
\emergencystretch=4em

\geometry{paperwidth=215mm,paperheight=290mm,hmargin=15mm,vmargin=16mm}
\begin{document}
\section*{P0067 Current delivery routes and future driver dispatch}
\subsection*{Mathematical formulation}
\begin{align}
a_{r,i_1}&=d_{0i_1}^n,\nonumber\\
a_{r,i_\ell}&=a_{r,i_{\ell-1}}+s+d_{i_{\ell-1}i_\ell}^n,
&&\ell=2,\ldots,h,\tag{D1}\\
\ell_r&=a_{r,i_h}+s+d_{i_h0}^n.\tag{D2}
\end{align}

\begin{equation}
\sum_{\ell=1}^h q_{i_\ell}^n\le Q,\qquad
 a_{r,i_\ell}\le L_{\max}\quad(\ell=1,\ldots,h).\tag{D3}
\end{equation}

\begin{equation}
c_r=\sum_{\ell=1}^h a_{r,i_\ell},\qquad
\alpha_{i,r}=\mathbf1\{i\in r\},\qquad
\beta_{p,r}=\mathbf1\{\ell_r>t_p-t_n\},\quad p>n.\tag{D4}
\end{equation}

\begin{align}
\min\quad&\sum_{r\in\mathcal R^n}c_r\theta_r+
\sum_{p=n+1}^N\sum_{k\in\mathcal K_p}h_p(k)x_k^p\tag{14}\\
\text{s.t.}\quad&
\sum_{m=n+1}^{p}\sum_{k=0}^{\bar K}\omega_m^p(k)x_k^m
\le\bar K-\zeta_p^n-\sum_{k=0}^{\bar K}kz_k^p,
&&p=n+1,\ldots,N,\tag{15}\\
&\sum_{k=0}^{\bar K}x_k^p=1,
&&p=n+1,\ldots,N,\tag{16}\\
&\sum_{k=0}^{\bar K}z_k^p=1,
&&p=n+1,\ldots,N,\tag{17}\\
&\sum_{r\in\mathcal R^n}\beta_{p,r}\theta_r
\le\sum_{k=0}^{\bar K}kz_k^p,
&&p=n+1,\ldots,N,\tag{18}\\
&\sum_{r\in\mathcal R^n}\theta_r\le\bar K-\zeta_n^n,\tag{19}\\
&\sum_{r\in\mathcal R^n}\alpha_{i,r}\theta_r=1,
&&i\in\mathcal I^n,\tag{20}\\
&x_k^p,z_k^p\in\{0,1\},
&&p=n+1,\ldots,N,\ k=0,\ldots,\bar K,\tag{21}\\
&\theta_r\in\{0,1\},&&r\in\mathcal R^n,\tag{22}\\
&x_k^p=0,&&p>n,\ k\notin\mathcal K_p.\tag{finite}
\end{align}

\subsection*{Definitions and model scope}
A delivery store is at a given dispatch time in a sequence of dispatch periods. All current orders are ready, their customer locations and strictly positive whole-number demands are known, and every demand fits one vehicle. The store has a given number of identical drivers and vehicles of a given capacity. For each current or later dispatch time it already knows how many drivers sent earlier will still be occupied. Distinct-location travel times, a common customer service duration and a maximum time from readiness to delivery are given. The store chooses current routes and plans future dispatch counts using given estimated future delivery costs and estimated driver occupation from future dispatches. Those estimates describe an approximate operating plan, and its best value is not claimed to be the exact best value of the underlying stochastic dynamic service.

Every available current route starts at the store now, visits a nonempty sequence of distinct current customers and returns to the store without an intermediate return for redispatch. The first delivery's elapsed time is travel time from the store; each later delivery's elapsed time adds the previous customer's service duration and the travel time from that customer. Return duration adds the final customer's service duration and travel back to the store. Total demand on the route cannot exceed vehicle capacity, and each delivery must occur within the delivery-time limit. Return itself need not meet that delivery limit. All routes meeting these requirements are eligible, not just a selected shortlist. A route's cost is the sum of its customer delivery elapsed times. It continues occupying a driver at a later dispatch time only if its return duration is strictly greater than the elapsed time to that dispatch; returning exactly then makes the driver available.

Each current route has a yes-or-no route dispatch record $\theta_r$. Each future period selects exactly one planned dispatch count between zero and total drivers using yes-or-no dispatch-count records $x_k^p$. Counts whose given future delivery cost is infinite are unavailable and their records are fixed off; only finite costs enter the total. Each future period also selects exactly one whole-number reservation between zero and total drivers using yes-or-no reservation records $z_k^p$. The reservation is at least the number of selected current routes still away at that future dispatch. It may exceed the actual number still away and is not required to equal it.

The number of current selected routes cannot exceed total drivers less the drivers already occupied at the current time. Each current customer belongs to exactly one selected route, with no splitting or postponement. At each future dispatch time, add the given estimated still-occupied drivers from every future dispatch between the next period and that time, using the count selected for each originating dispatch. This total cannot exceed fleet size less the already-known occupation from earlier dispatches and less the reservation for current routes at that future time. An originating future dispatch uses its selected driver count immediately; selecting zero contributes zero occupation. The estimates may otherwise be nonnegative real amounts.

The store seeks the lowest total of selected current route delivery-time costs and the supplied future delivery cost for each selected future dispatch count. No extra current-period estimate or vehicle penalty is added. At the last period there are no future choices or future-cost terms. If no current orders exist, no nonempty current route exists and no driver need be dispatched now. Future zero dispatch remains available precisely when its supplied delivery cost is finite.

\subsection*{Inline notation and relations}
$t_n$; $1,\ldots,N$; $t_1<\cdots<t_N$; $\mathcal I^n$; $0$; $i\in\mathcal I^n$; $q_i^n\in\mathbb Z_{>0}$; $q_i^n\le Q$; $Q>0$; $\bar K$; $\zeta_p^n\in\{0,\ldots,\bar K\}$; $t_p$; $p=n,\ldots,N$; $d_{ij}^n\ge0$; $i,j\in\{0\}\cup\mathcal I^n$; $t_{ij}^n$; $s\ge0$; $t_p=0$; $L_{\max}$; $p=n+1,\ldots,N$; $k=0,\ldots,\bar K$; $h_p(k)$; $\mathbb E_{\mathbb Q^p}[\mathcal H^s(k,\mathcal I^p,\mathbf q^p)]$; $\mathbb Q^p$; $+\infty$; $\mathcal K_p=\{k:h_p(k)<+\infty\}$; $x_k^p=0$; $\omega_m^p(k)\ge0$; $m$; $k$; $p\ge m$; $\omega_m^m(k)=k$; $\omega_m^p(0)=0$; $\mathcal R^n$; $r=(0,i_1,\ldots,i_h,0)$; $i_1,\ldots,i_h\in\mathcal I^n$; $\ell$; $\theta_r\in\{0,1\}$; $r$; $x_k^p\in\{0,1\}$; $p$; $z_k^p\in\{0,1\}$; $\zeta$; $\zeta_n^{p}$; $\zeta_p^n$; $n=N$; $\mathcal I^n=\varnothing$.
\end{document}